\documentclass[crop,tikz,border=5pt,convert={outext=.svg,command=\unexpanded{pdf2svg \infile\space\outfile}},multi=false]{standalone}[2012/04/13]
\makeatletter


\usetikzlibrary{calc}
\usepackage{relsize}
\tikzset{fontscale/.style = {font=\relsize{#1}}}


\definecolor{hasklogo@commodoreblue}{HTML}{1B3274}
\definecolor{hasklogo@commodorered}{HTML}{DB1B1B}
\pgfkeys{/hasklogo/commodore/.style={
  lambdabehind,
  logo   color=hasklogo@commodoreblue,
  lambda color=hasklogo@commodorered
}}
%
% The grayscale style likes the lambda in front and is, well, gray.
%
\pgfkeys{/hasklogo/grayscale/.style={
  lambdainfront,
  logo   color=black!40,
  lambda color=black!60
}}
%
% The seventies style likes the lambda behind and comes in seventies brown/orange combinations :)
%
\pgfkeys{/hasklogo/seventies/.style={
  lambdabehind,
  logo   color=orange,
  lambda color=brown!80!black
}}

\newif\ifhasklogo@lambdasplit@\hasklogo@lambdasplit@false
\pgfkeys{/hasklogo/first  rangle color/.store in=\hasklogo@first@rangle@color}
\pgfkeys{/hasklogo/second rangle color/.store in=\hasklogo@second@rangle@color}
\pgfkeys{/hasklogo/lambda        color/.store in=\hasklogo@lambda@color}
\pgfkeys{/hasklogo/lower  equals color/.store in=\hasklogo@lower@equals@color}
\pgfkeys{/hasklogo/upper  equals color/.store in=\hasklogo@upper@equals@color}
\pgfkeys{/hasklogo/xunit/.style={/tikz/x=#1}}
\pgfkeys{/hasklogo/yunit/.style={/tikz/y=#1}}
\pgfkeys{/hasklogo/unit/.style={/tikz/x=#1,/tikz/y=#1}}
\pgfkeys{/hasklogo/outlines/.style={fill=none,draw}}
\pgfkeys{/hasklogo/logo color/.style={
  first  rangle color=#1,
  second rangle color=#1,
  lambda        color=#1,
  lower  equals color=#1,
  upper  equals color=#1
}}
\pgfkeys{/hasklogo/lambdabehind/.code={\global\hasklogo@lambdasplit@true}}
\pgfkeys{/hasklogo/lambdainfront/.code={\global\hasklogo@lambdasplit@false}}
\pgfkeys{/hasklogo/.unknown/.code={%
  \let\searchname=\pgfkeyscurrentname
  \pgfkeysalso{\searchname/.try=#1,/tikz/\searchname/.retry=#1,/pgf/\searchname/.retry=#1}%
}}


\begin{document}
\begin{tikzpicture}[x=1ex,y=1ex]
\coordinate(hasklogoref);

  \tikzstyle{every path} = [fill,/hasklogo/.cd,lambdainfront,grayscale,seventies,lambdainfront,xunit=1,yunit=1];
  \path[color=\hasklogo@first@rangle@color] (hasklogoref)
    -- ++(4,6) -- ++(-4,6) -- ++(3,0) -- ++(4,-6) -- ++(-4,-6) -- cycle;
  \ifhasklogo@lambdasplit@
    \path[color=\hasklogo@lambda@color] (hasklogoref) ++(11,6)
    -- ++(4,-6) -- ++(-3,0) -- ++(-2.5,3.75) -- cycle;
  \path[color=\hasklogo@second@rangle@color] (hasklogoref) ++(4,0)
    -- ++(4,6) -- ++(-4,6) -- ++(3,0) -- ++(4,-6) -- ++(-4,-6) -- cycle;
  \else
    \path[color=\hasklogo@lambda@color] (hasklogoref) ++(8,6)
    -- ++(-4,6) -- ++(3,0) -- ++(8,-12) -- ++(-3,0) -- ++(-2.5,3.75) -- ++(-2.5,-3.75) -- ++(-3,0) -- cycle;
  \fi
    \path[fill=none,draw=none] (hasklogoref) -- coordinate[pos=1](c0) ++(8,12) -- coordinate[pos=1](c1) ++(8,-12);
  \foreach \miny/\maxy/\part in {3.5/5.5/lower,6.5/8.5/upper} {
    \path[fill=none,draw=none]
      (hasklogoref) -- coordinate[pos=1](hsleq0) ++(17,\miny)
      (hasklogoref) -- coordinate[pos=1](hsleq1) ++(17,\maxy);
    \coordinate(s0) at (intersection cs:
        first line={(hsleq0)--(hsleq0 -| c0)},
        second line={(c0)--(c1)}
        );
    \coordinate(s1) at (intersection cs:
        first line={(hsleq1)--(hsleq1 -| c0)},
        second line={(c0)--(c1)}
        );
    \path[color=\csname hasklogo@\part @equals@color\endcsname] (s0) -- (s1) -| (hsleq0) -- cycle;
  }

  \node at ($ (hasklogoref) + (30,5.5) $) [fontscale=2] {\TeX};
\end{tikzpicture}
\end{document}
